\chapter{Conditions}\label{ch:conditions}

A \term{condition} is an object that represents an exceptional
situation: an error, a warning, or a request to recover in some
particular way. A program \term{signals} a condition when the situation
arises, and the condition system finds a \term{handler} for it, a
function that some caller established in advance, and calls it. This
chapter follows chapter~13 of the 1992 manual (\dylanbook{chapter 13}),
which both implementations provide complete; it also describes what each
implementation does with a condition no handler takes, and the errors
the implementations themselves signal. As in chapter~\ref{ch:numbers}, a
class in a result is written by its name alone.

\section{How the condition system works}\label{sec:conditions-overview}

\index{signal}\index{handler}
Dylan's conditions are objects, instances of subclasses of
\code{<condition>}, and a handler is chosen by the condition's class: a
handler for \code{<error>} applies to every error, one for
\code{<type-error>} only to type errors. Handlers are established
dynamically, for the duration of a form's evaluation, and the most
recently established applicable handler is tried first.

Signalling is a function call. When a condition is signalled, the
signaller is still active and the handler runs ``on top of'' it. The
handler then does one of three things:

\begin{itemize}
\item It \textbf{exits} non-locally, abandoning the computation that
  signalled, to a point established outside it. This is the common case,
  and \code{handler-case} does it for you.
\item It \textbf{returns} values, which become the values of
  \code{signal} in the signaller. This is allowed only for conditions
  whose \term{recovery protocol} says what the values mean, such as
  warnings.
\item It \textbf{declines}, by calling the next handler, letting an
  older handler decide.
\end{itemize}

\noindent If every handler declines, or there are none, the condition's
\term{default handler} runs. For a warning it prints the message and
returns; for an error it stops the program, or in the Common Lisp
implementation enters the debugger.

Recovery can also be formal. A \term{restart} is a condition that a
handler signals to ask for a particular recovery, and a \term{restart
handler}, established by the code that knows how to recover, handles it.
Section~\ref{sec:conditions-restarts} describes restarts.

\section{The condition classes}\label{sec:conditions-classes}

\index{condition classes}
The predefined condition classes are these:

\begin{itemize}[nosep]
\item \code{<condition>}\index{<condition>}
  \begin{itemize}[nosep]
  \item \code{<serious-condition>}\index{<serious-condition>}
    \begin{itemize}[nosep]
    \item \code{<error>}\index{<error>}
      \begin{itemize}[nosep]
      \item \code{<simple-error>}\index{<simple-error>}
      \item \code{<type-error>}\index{<type-error>}
      \end{itemize}
    \end{itemize}
  \item \code{<warning>}\index{<warning>}
    \begin{itemize}[nosep]
    \item \code{<simple-warning>}\index{<simple-warning>}
    \end{itemize}
  \item \code{<restart>}\index{<restart>}
    \begin{itemize}[nosep]
    \item \code{<simple-restart>}\index{<simple-restart>}
    \item \code{<abort>}\index{<abort>}
    \end{itemize}
  \end{itemize}
\end{itemize}

\begin{description}
\item[\code{<condition>}] The class of all conditions. Its default
  handler returns \code{\#f}, so a plain condition nobody handles is
  ignored.
\item[\code{<serious-condition>}] Conditions that cannot safely be
  ignored. Its default handler is the debugger of last resort
  (section~\ref{sec:conditions-unhandled}).
\item[\code{<error>}] Serious conditions that represent something wrong
  with the program, as distinct from its environment; a handler for
  \code{<error>} does not catch other serious conditions.
\item[\code{<simple-error>}] An error that is just a message: a format
  string and arguments. They are given by the init-keywords
  \code{format-string:} and \code{format-arguments:}, and the functions
  \code{condition-format-string}\index{condition-format-string} and
  \code{condition-format-arguments}\index{condition-format-arguments}
  read them. \code{error} makes one from a string.
\item[\code{<type-error>}] An error saying a value was not of the
  expected type. Its init-keywords are \code{value:} and \code{type:},
  read by \code{type-error-value}\index{type-error-value} and
  \code{type-error-expected-type}\index{type-error-expected-type}.
  \code{check-type} signals one, and so do the implementations when a
  typed slot or variable is given a value of the wrong type.
\item[\code{<warning>}] Conditions worth reporting but not stopping
  for. The default handler prints the message and returns \code{\#f}.
\item[\code{<simple-warning>}] A warning that is a format string and
  arguments, like \code{<simple-error>}; \code{signal} makes one from a
  string.
\item[\code{<restart>}] Requests to recover. The default handler signals
  an error saying no handler was established for the restart. The
  init-keyword \code{condition:} is accepted, so that a restart can name
  the condition it recovers from.
\item[\code{<simple-restart>}] The restart \code{cerror} and
  \code{break} establish handlers for.
\item[\code{<abort>}] The restart that asks to abandon the current
  command and return to something like a command loop.
\end{description}

\noindent Conditions are made with \code{make} like any object. A
condition prints as its message:

\begin{dylan}
(print (make <simple-error> format-string: "bad value: ~s"
                            format-arguments: '(42)))
; prints bad value: 42
\end{dylan}

\section{Signalling}\label{sec:conditions-signalling}

\index{error}
\code{(error \meta{condition})} signals a condition and never returns.
\code{(error \meta{string} \meta{argument} ...)} makes a
\code{<simple-error>} from a format string and arguments and signals it.
This is the usual way to report an error:

\begin{dylan}
(define-method withdraw (balance amount)
  (if (> amount balance)
      (error "cannot withdraw ~a from a balance of ~a" amount balance)
      (- balance amount)))
\end{dylan}

\noindent If a handler for the error returns instead of exiting,
\code{error} does not return either: it reports the error as unhandled
(section~\ref{sec:conditions-unhandled}).

\index{signal}
\code{(signal \meta{condition})} signals a condition and answers whatever
values the handler returns. \code{(signal \meta{string} \meta{argument}
...)} makes a \code{<simple-warning>} and signals it. With no handler, a
warning is printed and \code{signal} answers \code{\#f}:

\begin{dylan}
(signal "disk is ~a% full" 91)
; prints a warning:  disk is 91% full
; => #f
\end{dylan}

\noindent The compiler to C prints the warning on standard error as
\code{Dylan warning: disk is 91\% full}; the Common Lisp implementation
uses Common Lisp's \code{warn}, which prints
\code{WARNING: disk is 91\% full}.

\index{format directives}
The format strings of \code{error}, \code{signal}, \code{cerror} and
\code{break} understand the directives \code{\textasciitilde{}a} and
\code{\textasciitilde{}d} (the argument as \code{display} prints it),
\code{\textasciitilde{}s} (as \code{print} prints it),
\code{\textasciitilde{}\%} (a new line) and \code{\textasciitilde{}\textasciitilde{}}
(a tilde). Other directives are copied as they are, and arguments no
directive uses are printed after the message, separated by spaces.

\index{cerror}
\code{(cerror \meta{restart-description} \meta{condition-or-string}
\meta{argument} ...)} is a \term{continuable error}. It signals the error
as \code{error} does, but first establishes a handler for
\code{<simple-restart>} described by \meta{restart-description}. If a
handler (or, in the Common Lisp implementation, a person at the
debugger) signals a \code{<simple-restart>}, \code{cerror} returns
\code{\#f} and the program carries on with the correction the
description promised:

\begin{dylan}
(define-method careful-sqrt (x)
  (if (negative? x)
      (begin (cerror "Use the absolute value" "negative argument ~a" x)
             (sqrt (abs x)))
      (sqrt x)))

(handler-bind (<error> (method (c next) (signal (make <simple-restart>))))
  (careful-sqrt -16))                            ; => 4.0
\end{dylan}

\index{check-type}
\code{(check-type \meta{value} \meta{type})} answers the value if it is
an instance of the type, and signals a \code{<type-error>} if not:

\begin{dylan}
(check-type 5 <integer>)        ; => 5
(handler-case (check-type "abc" <integer>)
  ((<type-error> condition: c)
   (list (type-error-value c) (type-error-expected-type c))))
; => ("abc" <integer>)
\end{dylan}

\noindent The implementations differ when a handler returns from the
type error: the compiler to C signals it with \code{error}, so a
returning handler makes it unhandled; the Common Lisp implementation
signals it with \code{signal}, so \code{check-type} answers whatever the
handler returned.

\index{abort}
\code{(abort)} signals an \code{<abort>} with \code{error}; it never
returns. A handler for \code{<abort>} decides what abandoning means:

\begin{dylan}
(handler-case (begin (abort) 'not-reached)
  ((<abort>) 'aborted))         ; => aborted
\end{dylan}

\index{break}
\code{(break)} or \code{(break \meta{string} \meta{argument} ...)} is a
breakpoint. In the Common Lisp implementation it enters the debugger
with the message, offering a restart that continues; in the compiler to
C, which has no debugger, it prints \code{Dylan break:} and the message
on standard error and continues. Either way it then answers \code{\#f}.

\section{Handling with handler-case}\label{sec:conditions-handler-case}

\index{handler-case}
\code{handler-case} is the usual way to handle conditions: it evaluates a
form, and if a condition of one of the listed classes is signalled during
it, abandons the form and evaluates the matching clause instead.

\begin{dylan}
(handler-case protected-form
  ((type option ...) form ...)
  ...)
\end{dylan}

\noindent If the protected form finishes, \code{handler-case} answers its
values. If a condition is signalled that is an instance of a clause's
\meta{type} (a class, or a singleton), the protected form is exited and
that clause's forms are evaluated, their last value being the answer, or
\code{\#f} if there are none. The clauses are tried in order, so put the
more specific classes first. The options are:

\begin{description}
\item[\code{condition: \meta{name}}] binds the variable \meta{name} to the
  condition in the clause's forms;
\item[\code{test: \meta{function}}] a predicate the condition must
  also satisfy for the clause to apply;
\item[\code{description: \meta{string}}] a description of the handler,
  which introspection (\code{do-handlers}) can see.
\end{description}

\begin{dylan}
(handler-case (error "bad value: ~s" 42)
  ((<simple-error> condition: c)
   (list (condition-format-string c) (condition-format-arguments c))))
; => ("bad value: ~s" (42))

(handler-case (error "code ~a" 7)
  ((<simple-error> test: (method (c) (= (first (condition-format-arguments c)) 8)))
   'eight)
  ((<error>) 'other))
; => other

(handler-case (error "x")
  ((<error>) 'first)
  ((<simple-error>) 'second))
; => first

(handler-case (values 1 2)
  ((<error>) 'oops))
; => 1 and 2
\end{dylan}

\noindent The clause's forms run after the exit, in the dynamic
environment of the \code{handler-case} itself, so its handlers are no
longer in force there. A singleton type handles one particular
condition object:

\begin{dylan}
(define c1 (make <simple-error> format-string: "one"))
(define c2 (make <simple-error> format-string: "two"))
(define-method try (c)
  (handler-case (error c)
    (((singleton c1)) 'got-c1)
    ((<simple-error> condition: e) (list 'other (condition-format-string e)))))
(list (try c1) (try c2))     ; => (got-c1 (other "two"))
\end{dylan}

\section{Handling with handler-bind}\label{sec:conditions-handler-bind}

\index{handler-bind}\index{next-handler}
\code{handler-bind} is the primitive form, and both implementations
have it. It establishes one handler, a function, while its body is
evaluated, and answers the body's values:

\begin{dylan}
(handler-bind (type function option ...)
  form ...)
\end{dylan}

\noindent The type, function and options are evaluated before the body;
the options are \code{test:} and \code{description:}, as for
\code{handler-case}. When an applicable condition is signalled, the
function is called with two arguments, the condition and a
\term{next-handler} function, without exiting anything. It may exit
non-locally (with \code{bind-exit}, chapter~\ref{ch:control}); it may
return values, which \code{signal} then answers; or it may decline by
calling the next-handler with no arguments, which tries the next older
applicable handler.

\begin{dylan}
(handler-bind (<warning> (method (c next) (print 'outer) 'outer-value))
  (handler-bind (<warning> (method (c next) (print 'inner-declines) (next)))
    (signal "low battery")))
; prints inner-declines, then outer
; => outer-value
\end{dylan}

\noindent There is no ``condition wall'': while a handler runs, every
handler established before it is still in force, including those
established inside the signalling code. That is what lets a handler
signal a restart that the signaller handles
(section~\ref{sec:conditions-restarts}). It also means that a handler
which signals a condition of its own class, or causes an error it also
handles, can recurse until the stack is exhausted.

\section{Default handlers}\label{sec:conditions-default-handler}

\index{default-handler}
When no dynamic handler handles a condition, \code{signal} calls
\code{default-handler}, a generic function, on it. Its predefined methods
are:

\begin{center}
\begin{tabular}{ll}
\toprule
Method on & Does \\
\midrule
\code{<condition>} & answers \code{\#f} \\
\code{<serious-condition>} & reports the condition as unhandled \\
\code{<warning>} & prints the message, answers \code{\#f} \\
\code{<restart>} & signals an error: no handler was established for the restart \\
\bottomrule
\end{tabular}
\end{center}

\noindent A condition class of your own inherits the method of its
superclass. A condition that is neither serious nor a warning is simply
ignored when unhandled:

\begin{dylan}
(define-class <note> (<condition>)
  (text init-keyword: text:))

(signal (make <note> text: "hello"))          ; => #f
(handler-case (signal (make <note> text: "hi"))
  ((<note> condition: n) (text n)))           ; => "hi"
\end{dylan}

\noindent A program may call \code{default-handler} itself, and may add
methods to it for its own classes.

\section{Restarts}\label{sec:conditions-restarts}

\index{restart}
A restart is a condition that asks for a particular way of recovering.
The code that knows how to recover establishes a handler for a restart
class, usually with \code{handler-case}, around the code that may fail;
a handler for the original condition, which knows which recovery it
wants, signals the restart. Values a recovery needs travel in the
restart's slots.

This program parses digits. When a character is not a digit it signals
an error, but offers a \code{<use-value>} restart that supplies a value
instead; the caller's handler chooses that recovery:

\begin{dylan}
(define-class <use-value> (<restart>)
  (value init-keyword: value:))

(define-method parse-digit ((c <character>))
  (if (and (>= c #\0) (<= c #\9))
      (- (as <integer> c) (as <integer> #\0))
      (handler-case (error "not a digit: ~s" c)
        ((<use-value> condition: r) (value r)))))

(handler-bind (<simple-error>
               (method (c next) (signal (make <use-value> value: -1))))
  (map parse-digit '(#\4 #\x #\2)))
; => (4 -1 2)
\end{dylan}

\noindent The handler for \code{<simple-error>} runs while
\code{parse-digit}'s \code{handler-case} is still in force, so its
\code{signal} reaches the restart handler, which exits from the
\code{handler-case} with the value.

\index{<simple-restart>}
\code{cerror} establishes a \code{<simple-restart>} handler, as shown in
section~\ref{sec:conditions-signalling}; signalling a
\code{<simple-restart>} where none is established is an error. The
\code{condition:} init-keyword lets a restart be tied to one condition,
so that a \code{test:} on its handler can accept only restarts meant
for that condition (\dylanbook{chapter 13} has an example).

\section{Introspection}\label{sec:conditions-introspection}

\index{do-handlers}
\code{(do-handlers \meta{function})} calls the function on each handler
in force, the most recent first, with four arguments: the handler's type,
test, function and description. An absent test is given as a function
that answers \code{\#t}, and an absent description as \code{\#f}.

\begin{dylan}
(handler-bind (<warning> (method (c n) #f) description: "ignore warnings")
  (handler-bind (<error> (method (c n) #f))
    (do-handlers (method (type test function description)
                   (print (list type description))))))
; prints (<error> #f), then (<warning> "ignore warnings")
\end{dylan}

\noindent At the Common Lisp implementation's prompt the list ends with
the loop's own handler, an \code{<abort>} described as
\code{"Return to the read-eval-print loop."}

\index{return-allowed?}\index{return-description}
\index{restart-query}\index{return-query}
Four generic functions let an interactive debugger ask whether a
condition may be handled by returning (\code{return-allowed?}), what
returning would mean (\code{return-description}), and the user for the
contents of a restart (\code{restart-query}) or for values to return
(\code{return-query}). The compiler to C, having no debugger, defines
only methods that do nothing: all four answer \code{\#f}. In the Common
Lisp implementation \code{return-allowed?} answers \code{\#t} for a
warning, whose recovery protocol allows returning, and
\code{return-description} then answers the string
\code{"Return; the value is ignored."}; its \code{return-query} reads
values typed by the user.

\section{Defining your own conditions}\label{sec:conditions-defining}

\index{condition!defining}
Define a condition class as a subclass of the predefined class closest
in meaning: \code{<error>} for a program error, \code{<warning>} for
something worth reporting, \code{<restart>} for a recovery. Slots carry
whatever a handler needs to know.

The simplest way to give a condition a message that both
implementations report is to subclass \code{<simple-error>} (or
\code{<simple-warning>}) and pass \code{format-string:} and
\code{format-arguments:} when making it:

\begin{dylan}
(define-class <bad-reading> (<simple-error>)
  (sensor init-keyword: sensor:))

(define-method report-bad (sensor value)
  (error (make <bad-reading> sensor: sensor
               format-string: "sensor ~a read ~a"
               format-arguments: (list sensor value))))

(handler-case (report-bad 'north 9999)
  ((<bad-reading> condition: c) (list (sensor c) c)))
; => (north sensor north read 9999)
\end{dylan}

\noindent Unhandled, it is reported as \code{sensor south read -1}. A
direct subclass of \code{<error>} with no message is reported by its
class only: in the compiler to C as
\code{Dylan error: \{instance of <oops>\}}, even if the class has a
\code{print-object} method (chapter~\ref{ch:printing}); that method is
used when the program prints the condition itself.

\section{Unhandled conditions}\label{sec:conditions-unhandled}

\index{unhandled condition}\index{debugger}
What happens to a serious condition that no handler takes, the
``debugger of last resort'' of the manual, is where the implementations
differ most.

\subsection{In the Common Lisp implementation}

An unhandled error enters the Common Lisp debugger, which shows the
message and offers restarts: Dylan's restart handlers in force, then
restarts that return to the Dylan prompt, then Common Lisp's own. Type
the number of a restart to choose it. Here a \code{cerror} offers its
restart:

\begin{repl}
? (define-method checked (n)
    (when (> n 10)
      (cerror "Use 10 instead" "~a is too big" n)
      (set! n 10))
    n)

;Value: checked

? (checked 99)

debugger invoked on a DYLAN::DYLAN-CONDITION in thread
#<THREAD tid=712328 "main thread" RUNNING {1200030003}>:
  Error: 99 is too big

Type HELP for debugger help, or (SB-EXT:EXIT) to exit from SBCL.

restarts (invokable by number or by possibly-abbreviated name):
  0: [DYLAN-RESTART] Use 10 instead
  1:                 Return to the read-eval-print loop.
  2: [DYLAN-REP    ] Return to the Dylan read-eval-print loop.
  3: [ABORT        ] Exit debugger, returning to top level.

((LAMBDA NIL :IN DYLAN::IMPLEMENTATION-SPECIFIC/INDUCE-ERROR))
0] 0

;Value: 10
\end{repl}

\noindent Choosing \code{Return to the read-eval-print loop} abandons the
form and prints \code{;Aborted}; \code{DYLAN-REP} abandons it
without the message. \code{break} offers a \code{CONTINUE}
restart, which makes \code{break} answer \code{\#f}.

An error signalled by Common Lisp itself (a Lisp type error inside a
built-in function, a floating-point overflow) is offered first to Dylan's
handlers as a \code{<simple-error>} whose message begins
\code{Lisp condition:}. If none takes it, it reaches the debugger as the
Lisp error it is, and the debugger offers \code{DYLAN-REP}, which
returns to the Dylan prompt, and Common Lisp's own restarts. In a program compiled to Lisp with
\code{dylan:compile-file} and loaded into the runtime, an unhandled error
likewise goes to the Common Lisp debugger (or, when SBCL runs with its
debugger disabled, ends SBCL with a backtrace). Chapter~\ref{ch:lisp} describes the debugger further.

\subsection{In the compiler to C}

A compiled program has no debugger. An unhandled serious condition in
the main thread prints \code{Dylan error:} and the condition's message
on standard error and ends the program with exit status 1. Anything the
program had written to standard output is flushed first.

\begin{shell}
$ cat range.dyl
(print "before")
(error "the value ~s is out of range" 42)
(print "after")
$ bin/dylanc range.dyl
$ ./range
"before"
Dylan error: the value 42 is out of range
$ echo $?
1
\end{shell}

\noindent The cleanup forms of \code{unwind-protect} that are in force are
\emph{not} run when a program stops this way. Some other cases:

\begin{center}
\begin{tabular}{p{0.38\textwidth}p{0.55\textwidth}}
\toprule
Program & Reported \\
\midrule
\code{(check-type "x" <integer>)} & \code{Dylan error: "x" is not an instance of \{the class <integer>\}} \\
a handler returns from \code{(error "nope")} & \code{Dylan error: a handler returned from error: nope} \\
\code{(abort)} with no handler & \code{Dylan error: no handler was established for the restart \{instance of <abort>\}} \\
\bottomrule
\end{tabular}
\end{center}

\noindent A condition a thread other than the main one does not handle
ends that thread instead, as chapter~\ref{ch:extensions} describes.

\section{Errors the implementations signal}\label{sec:conditions-runtime-errors}

\index{runtime errors}
Errors the implementations detect themselves are signalled as Dylan
conditions, so programs can handle them like any other. They are
\code{<type-error>}s when a value of the wrong type is stored in a typed
slot or variable, and \code{<simple-error>}s otherwise. In the compiler
to C:

\begin{center}
\begin{tabular}{p{0.36\textwidth}p{0.57\textwidth}}
\toprule
Cause & Message \\
\midrule
no method applies to the arguments & \code{no applicable method for f} \\
wrong number of arguments & \code{wrong number of arguments (2)} \\
a keyword the function does not take & \code{keyword argument not accepted: b:} \\
calling something not a function & \code{not a function: 5} \\
integer overflow & \code{integer overflow (integers are 63 bits)} \\
division by zero & \code{division by zero} \\
a typed slot or variable given a wrong value & \code{"s" is not an instance of \{the class <integer>\}} (a \code{<type-error>}) \\
reading an uninitialized slot & \code{the slot s is not initialized} \\
a variable with no value & \code{the variable x@dylan-user has no value} \\
\code{car} of a non-list & \code{not a list: 5} \\
an element that does not exist & \code{element: there is no element at the index 5} \\
\code{make} of an abstract class & \code{cannot make an instance of this class: \{the class <integer>\}} \\
\bottomrule
\end{tabular}
\end{center}

\noindent The compiler also warns, when compiling, about a variable that
is not defined anywhere (\code{warning: x is not defined}); the program
is still built, and signals the error above if the reference is
evaluated.

The Common Lisp implementation detects the same errors with messages of
its own, for example \code{generic-dispatch -- no applicable methods},
\code{invalid number of arguments: 3}, \code{illegal keyword argument b:
(a:)} and \code{unbound variable x}. An error that arises inside Common
Lisp, such as \code{(car 5)}, has a message beginning
\code{Lisp condition:} followed by Common Lisp's own description.

\begin{dylan}
(define-method f ((x <integer>)) x)

(handler-case (f "s")
  ((<error> condition: c) c))
; C:    no applicable method for f
; Lisp: generic-dispatch -- no applicable methods {the generic function f} ("s")

(define-class <account> (<object>)
  (balance type: <integer> init-value: 0))

(handler-case (set! (balance (make <account>)) "lots")
  ((<type-error> condition: c) (type-error-value c)))
; => "lots"
\end{dylan}
